\documentclass[11pt]{article}

\usepackage[margin=1in]{geometry}
\usepackage{booktabs}
\usepackage{hyperref}
\usepackage{url}
\usepackage{xcolor}

\hypersetup{
  colorlinks=true,
  linkcolor=blue,
  citecolor=blue,
  urlcolor=blue
}

\title{InfraBench: A Reproducible Benchmark for Evaluating AI Agents on Real Infrastructure Operations}

\author{
Thomas Chaigneau\\
Independent Researcher / Kubeply\\
\texttt{thomas@kubeply.com}
}

\date{Draft of May 5, 2026}

\begin{document}
\maketitle

\begin{abstract}
AI agents are increasingly evaluated on software engineering, desktop-use, and
terminal-use benchmarks, but infrastructure operations remain underrepresented
despite being a practical target for autonomous agents. We introduce
InfraBench, an open benchmark of realistic infrastructure tasks for evaluating
agents in reproducible execution environments. The first dataset,
\texttt{kubeply/kubernetes-core}, contains 58 Kubernetes operations tasks
covering workload health, service connectivity, scheduling, storage, access
control, migration, maintenance, and incident response. Each task includes a
natural-language operator instruction, an isolated environment, a deterministic
oracle solution, and an executable verifier. The benchmark is compatible with
Harbor and is designed to measure whether an agent can diagnose a broken
infrastructure state, make a minimal repair, and leave the system in a
verifiably correct condition. Preliminary runs across nine frontier and
open-model agent configurations show aggregate success rates between 68.9\%
and 84.5\%, with materially lower performance on hard tasks than on easy tasks.
These results suggest that infrastructure operations expose a different failure
surface from conventional coding benchmarks: agents often solve simple manifest
or routing defects, but struggle with multi-resource state preservation,
rollout coordination, and policy-sensitive repairs. We release the task
definitions, metadata, oracle scripts, and result normalization tooling under
Apache-2.0 to support reproducible evaluation and future infrastructure-agent
research.
\end{abstract}

\section{Introduction}

AI agents are moving from single-turn code generation toward long-horizon
computer work. Existing benchmarks have made this transition measurable in
software maintenance, desktop environments, and terminal workflows
\cite{jimenez2024swebench,xie2024osworld,merrill2026terminalbench}. These
benchmarks are valuable because they replace synthetic prompts with executable
tasks, hidden tests, and stateful environments. However, infrastructure and
platform engineering remain weakly covered as a benchmark domain.

This gap matters. Production infrastructure work differs from ordinary coding
tasks in several ways. The agent must reason about live or staged system state,
understand operational invariants, preserve identities and security boundaries,
coordinate rollouts, and avoid changes that appear to fix a symptom while
violating the intended platform contract. A Kubernetes incident can require
diagnosing a Deployment, Service, Ingress, Secret, NetworkPolicy, PodDisruption
Budget, and node condition together. These tasks are not purely code-editing
problems; they are state repair problems.

InfraBench is an open benchmark intended to make this class of work measurable.
Its first dataset, \texttt{kubeply/kubernetes-core}, focuses on Kubernetes
operator tasks. Each task presents the agent with a concise instruction and an
environment containing a broken infrastructure state. The agent must inspect
the environment, perform the repair, and pass an executable verifier. The
verifier checks the operator outcome rather than simply matching an expected
patch.

This paper makes three contributions:

\begin{enumerate}
  \item We define a benchmark format for realistic infrastructure operations
  tasks with natural-language instructions, isolated environments, oracle
  solutions, and deterministic verifiers.
  \item We release \texttt{kubeply/kubernetes-core}, a 58-task Kubernetes
  dataset spanning eight task categories and three difficulty tiers.
  \item We report preliminary multi-agent results and analyze how performance
  varies by task difficulty, showing that hard infrastructure tasks remain
  substantially less reliable than easy diagnosis tasks.
\end{enumerate}

\section{Related Work}

SWE-bench evaluates language models on real software engineering issues drawn
from GitHub repositories \cite{jimenez2024swebench}. It established that
repository-scale repair is a useful and difficult target for agent evaluation.
InfraBench follows the same broad principle of executable, realistic tasks, but
targets infrastructure operations rather than application code maintenance.

OSWorld evaluates multimodal agents in real desktop environments
\cite{xie2024osworld}. Its emphasis on stateful environments and executable
success criteria is similar to InfraBench, but the task domain is general
computer use. Terminal-Bench focuses on hard command-line tasks in terminal
environments and includes unique environments, human-written solutions, and
comprehensive tests \cite{merrill2026terminalbench}. InfraBench is closest in
spirit to Terminal-Bench, but narrows the domain to infrastructure operations
and adds task metadata around operational categories, Kubernetes concepts, and
difficulty calibration.

Kubernetes itself is a mature and widely used orchestration system
\cite{kubernetes2026}. It is a natural first target because infrastructure
agents frequently need to reason about Kubernetes resources, runtime state,
controller behavior, network policy, scheduling, storage identity, and
configuration propagation.

\section{Benchmark Design}

\subsection{Task Shape}

InfraBench tasks follow the Harbor-compatible task shape
\cite{harbor2026,infrabench2026}:

\begin{verbatim}
task-name/
|-- instruction.md
|-- task.toml
|-- environment/
|-- solution/
|   `-- solve.sh
`-- tests/
    `-- test.sh
\end{verbatim}

\texttt{instruction.md} is the only prompt shown to the agent. The
\texttt{environment/} directory contains the starting state. The
\texttt{solution/solve.sh} script is an oracle repair used to validate that the
task is solvable. The \texttt{tests/test.sh} verifier writes the reward result.

The benchmark intentionally separates three roles. The task prompt describes
the operator-facing problem. The environment contains the evidence and broken
state. The verifier encodes the hidden correctness criteria. This separation is
important because infrastructure tasks are vulnerable to prompt leakage: naming
the exact resource field or diagnosis in the prompt can turn a medium task into
a mechanical edit.

\subsection{Operator Realism}

Tasks are designed around one-sentence operator stories: a platform engineer
needs to restore a workload, route, policy, migration, or maintenance outcome
because a user-visible or operator-visible condition is broken. The agent is
usually expected to make a minimal repair rather than rewrite the environment.
Verifiers reject common bypasses such as deleting policies, replacing unrelated
resources, broadening access control, or mutating verifier-trusted baselines.

The current Kubernetes dataset uses neutral namespaces and resource names where
possible. Medium and hard tasks avoid telling the agent the exact root cause.
For example, a prompt may state that checkout records are failing without
naming the broken NetworkPolicy selector, stale Secret projection, or incorrect
headless Service identity that causes the failure.

\subsection{Scoring}

Each task returns a reward through the verifier. The current dataset uses
binary pass/fail scoring, and aggregate score is the mean reward over all tasks
in the selected dataset. Future versions may add partial credit where the
operator outcome naturally decomposes into independent sub-goals, but the first
release favors simple interpretability.

\section{The Kubernetes-Core Dataset}

\texttt{kubeply/kubernetes-core} contains 58 tasks. Table
\ref{tab:dataset-difficulty} shows the difficulty split, and Table
\ref{tab:dataset-category} shows the category split.

\begin{table}[h]
\centering
\begin{tabular}{lrr}
\toprule
Difficulty & Tasks & Share \\
\midrule
Easy & 22 & 37.9\% \\
Medium & 23 & 39.7\% \\
Hard & 13 & 22.4\% \\
\midrule
Total & 58 & 100.0\% \\
\bottomrule
\end{tabular}
\caption{Task counts by difficulty in \texttt{kubeply/kubernetes-core}.}
\label{tab:dataset-difficulty}
\end{table}

\begin{table}[h]
\centering
\begin{tabular}{lr}
\toprule
Category & Tasks \\
\midrule
Workload health & 11 \\
Service connectivity & 10 \\
Scheduling and capacity & 9 \\
Access and isolation & 8 \\
Platform APIs and controllers & 7 \\
Migration and maintenance & 6 \\
Storage and state & 4 \\
Configuration and secrets & 3 \\
\bottomrule
\end{tabular}
\caption{Task counts by operational category.}
\label{tab:dataset-category}
\end{table}

Most tasks are live-cluster debugging or incident-response scenarios. The
scenario-type distribution is 36 live-cluster debug tasks, 16 incident-response
tasks, two maintenance tasks, two migration tasks, one policy-authoring task,
and one upgrade-readiness task. This distribution reflects the first dataset's
goal: measure practical Kubernetes repair and diagnosis before expanding into
Terraform, observability, or cloud-provider operations.

\section{Preliminary Evaluation}

We ran the 58-task dataset through Harbor-compatible agent configurations. The
current results should be read as preliminary because the run artifacts have
not yet been attached to an immutable paper release. They are useful for
estimating whether the benchmark differentiates agents and difficulties.

Table \ref{tab:model-results} reports aggregate pass rates. The best observed
configuration solved 49 of 58 tasks. The lowest observed configuration solved
40 of 58 tasks. This range is narrow enough to suggest that many easy tasks are
within current agent capability, but wide enough to show meaningful
configuration and model differences.

\begin{table}[h]
\centering
\begin{tabular}{llrrr}
\toprule
Provider & Model / setting & Passed & Total & Score \\
\midrule
OpenAI & GPT-5.5 medium & 49 & 58 & 84.5\% \\
OpenAI & GPT-5.5 high & 46 & 58 & 79.3\% \\
OpenAI & GPT-5.4 high & 46 & 58 & 79.3\% \\
Anthropic & Claude Sonnet 4.6 high & 46 & 58 & 79.3\% \\
Anthropic & Claude Opus 4.6 high & 46 & 58 & 79.3\% \\
Google & Gemini 3 Flash Preview & 43 & 58 & 74.1\% \\
Moonshot & Kimi K2.5 & 41 & 58 & 70.7\% \\
Moonshot & Kimi K2.6 & 40 & 58 & 69.0\% \\
Z.ai & GLM-5.1 & 40 & 58 & 69.0\% \\
\bottomrule
\end{tabular}
\caption{Preliminary aggregate results on \texttt{kubeply/kubernetes-core}.}
\label{tab:model-results}
\end{table}

Table \ref{tab:difficulty-results} breaks the same runs down by difficulty.
The strongest pattern is that easy tasks are almost saturated, while hard tasks
remain substantially harder. For the best aggregate run, the agent solved 21 of
22 easy tasks, 20 of 23 medium tasks, and 8 of 13 hard tasks. Several other
configurations solved all or nearly all easy tasks but less than half of the
hard tasks.

\begin{table}[h]
\centering
\begin{tabular}{lrrr}
\toprule
Model / setting & Easy & Medium & Hard \\
\midrule
GPT-5.5 medium & 21/22 & 20/23 & 8/13 \\
GPT-5.5 high & 21/22 & 17/23 & 8/13 \\
GPT-5.4 high & 20/22 & 19/23 & 7/13 \\
Claude Sonnet 4.6 high & 22/22 & 19/23 & 5/13 \\
Claude Opus 4.6 high & 21/22 & 18/23 & 7/13 \\
Gemini 3 Flash Preview & 22/22 & 16/23 & 5/13 \\
Kimi K2.5 & 22/22 & 15/23 & 4/13 \\
Kimi K2.6 & 21/22 & 15/23 & 4/13 \\
GLM-5.1 & 21/22 & 16/23 & 3/13 \\
\bottomrule
\end{tabular}
\caption{Preliminary pass counts by difficulty.}
\label{tab:difficulty-results}
\end{table}

\section{Error Surface}

The preliminary results suggest three recurring failure modes. First, agents
often succeed on local manifest or selector repairs but struggle when the
correct fix requires preserving stateful identity across multiple resources.
Second, agents can produce repairs that restore immediate readiness while
weakening an isolation boundary, such as by broadening policy or access control
more than the task allows. Third, agents sometimes fail to coordinate
configuration changes with rollout behavior: the static resource becomes
correct, but the running workload does not consume the intended state.

These failure modes are operationally important. In production infrastructure,
a patch that restarts the wrong workload, deletes a security policy, or creates
a replacement resource with a new identity can pass a shallow health check
while causing a later incident. InfraBench verifiers therefore test both the
visible outcome and the surrounding invariants where possible.

\section{Reproducibility and Release Process}

InfraBench is released as an Apache-2.0 repository with task definitions,
oracle scripts, verifiers, dataset manifests, and result-normalization tooling
\cite{infrabench2026}. The intended reproducibility path is:

\begin{enumerate}
  \item select a tagged dataset release;
  \item run Harbor against a task or dataset with a specified agent and model;
  \item preserve raw job artifacts;
  \item normalize the run into public \texttt{run.json} and \texttt{results.json}
  summaries;
  \item report the exact InfraBench commit and model/harness metadata.
\end{enumerate}

This paper draft is based on repository commit \texttt{33ca0d2bfda2}. The
camera-ready version should replace this commit with an immutable dataset tag
and attach public run artifacts for every table.

\section{Limitations}

The first limitation is scope. The current dataset is Kubernetes-only. This is
intentional: Kubernetes provides a dense, realistic infrastructure domain, but
it does not cover Terraform, observability, cloud-provider APIs, CI/CD systems,
or incident collaboration workflows.

The second limitation is authoring bias. Tasks are written by the benchmark
maintainer and reflect one view of realistic platform work. Future releases
should include external task review from platform engineers and SREs.

The third limitation is contamination. Once tasks are public, model providers
and agents may indirectly train on them. InfraBench should therefore version
datasets, preserve historical results, and add fresh task releases over time.

The fourth limitation is the absence of a human baseline in the current draft.
Human operator timing and success rates would make the difficulty labels more
defensible and should be added before submission if feasible.

\section{Conclusion}

InfraBench targets a practical gap in AI-agent evaluation: realistic
infrastructure operations. The first dataset,
\texttt{kubeply/kubernetes-core}, packages 58 Kubernetes repair tasks with
operator prompts, isolated environments, oracle solutions, and executable
verifiers. Preliminary runs show that current agents handle many easy
Kubernetes repairs but remain less reliable on hard tasks involving state
preservation, rollout coordination, and policy-sensitive changes. These results
support the case for infrastructure-specific benchmarks as a complement to
software engineering, desktop, and terminal benchmarks.

\section*{Acknowledgements}

The benchmark uses the Harbor task format and evaluation workflow. The author
thanks the open-source infrastructure and AI-agent communities for the tools
and prior benchmarks that shaped this work.

\bibliographystyle{plain}
\bibliography{references}

\end{document}
